\chapter{Architecture at a Glance}
\label{chap:architecture-at-a-glance}

\section{Introduction}
This chapter provides a high-level, visual, and quantitative summary of the Nexus PM software architecture. It highlights the technology choices, core architectural principles, cloud deployment targets, and structural metrics that define the platform. 

This section serves as a fast-track reference for technical architects, engineering managers, and reviewers before diving into the detailed functional and logical designs of the subsequent chapters.

\section{System Architecture Overview}
Nexus PM is built on a decoupled, cloud-native, and client-centric architecture. The client single-page application (SPA) executes in the user's browser, managing state transitions and interface interactions locally. The backend serves as a stateless API gateway, validating requests, enforcing access control rules, and executing business logic asynchronously where possible to minimize response latency. 

The runtime boundary is organized into distinct logical tiers:
\begin{itemize}
    \item \textbf{Presentation Tier:} React 19 single-page application distributed globally via Vercel Edge CDN nodes.
    \item \textbf{Gateway \& Logic Tier:} Stateless FastAPI ASGI service running containerized workloads on Render.
    \item \textbf{Persistence Tier:} Managed PostgreSQL (Supabase) handling ACID transaction boundaries, and Cloudflare R2 object storage managing file binary assets.
    \item \textbf{Caching \& Orchestration Tier:} Upstash serverless Redis handling transient data (sessions, rate limits, dashboard metrics caching).
    \item \textbf{Intelligence Tier:} Google Gemini 2.5 Flash API integrating natural language processing capabilities.
\end{itemize}

\section{Technology Stack Summary}
Table~\ref{tab:at_a_glance_tech_stack} summarizes the core technology selections, layers, and operational justifications for the Nexus PM ecosystem.

\begin{table}[h!]
\centering
\small
\renewcommand{\arraystretch}{1.4}
\begin{tabularx}{\textwidth}{l l X}
    \toprule[1.5pt]
    \rowcolor{primary!5}
    \textbf{\color{primary}Layer} & \textbf{\color{primary}Technology} & \textbf{\color{primary}Primary Architectural Rationale} \\
    \midrule
    Client SPA & React 19, TS, Vite & Fast component-based rendering, static type safety, and fast bundle compilation. \\
    State & Redux Toolkit, React Query & Decoupled client-side UI states and caching for backend queries. \\
    Styling & Tailwind CSS & Rapid interface styling via utility class constraints. \\
    API Service & FastAPI (Python 3.12) & Async-first performance, native OpenAPI docs, and type verification. \\
    ORM & SQLAlchemy, Alembic & Relational mapping, schema version control, and query abstraction. \\
    Database & PostgreSQL (Supabase) & ACID transactional safety, managed database maintenance, and backups. \\
    In-Memory Cache & Upstash Redis & Serverless rate limiting, session caching, and low latency reads. \\
    Object Storage & Cloudflare R2 & Zero data egress bandwidth billing, S3-compatible API. \\
    AI Engine & Gemini 2.5 Flash & High context reasoning limits, structured JSON outputs, low cost. \\
    Email Pipeline & Resend API & Developer-friendly transactional email delivery. \\
    Virtualization & Docker & Identical local and production runtime environments. \\
    \bottomrule[1.5pt]
\end{tabularx}
\caption{Nexus PM Core Technology Selection Summary}
\label{tab:at_a_glance_tech_stack}
\end{table}

\section{Architectural Principles}
The platform's design conforms to four core software engineering paradigms:
\begin{itemize}
    \item \textbf{Complete Client-Server Decoupling:} The frontend client application and backend API service communicate strictly over HTTPS endpoints, enabling independent development, deployment, and scalability.
    \item \textbf{Stateless Gateway Operation:} The API servers store no persistent state on local disk or local memory. All user sessions are validated statelessly using cryptographically signed JSON Web Tokens (JWT) or transient Redis tokens.
    \item \textbf{Non-Blocking Asynchronous Operations:} Heavy computational work, transactional email dispatches, and generative AI model requests are executed asynchronously outside of the primary HTTP request-response cycle.
    \item \textbf{Tenant Resource Isolation:} Multi-tenant isolation is enforced at the API controller layer through database query filters and role validation checks.
\end{itemize}

\pagebreak
\section{Key Architectural Figures}
The detailed structures, interactions, and schemas described throughout this document are represented by six primary diagrams:
\begin{itemize}
    \item \textbf{System Context \& Integration Topology:} Figure~\ref{fig:system_architecture} in Chapter~\ref{chap:high-level-design} (Page~\pageref{fig:system_architecture}) shows the operational boundaries and integrations.
    \item \textbf{Runtime Deployment Architecture:} Figure~\ref{fig:deployment_topology} in Chapter~\ref{chap:deployment-design} (Page~\pageref{fig:deployment_topology}) outlines container network routing.
    \item \textbf{Relational Database Schema (ERD):} Figure~\ref{fig:database_erd} in Chapter~\ref{chap:database-design} (Page~\pageref{fig:database_erd}) maps tables, primary/foreign keys, and data relationships.
    \item \textbf{Authentication and Identity Verification Flow:} Figure~\ref{fig:auth_flow_diagram} in Chapter~\ref{chap:security-design} (Page~\pageref{fig:auth_flow_diagram}) details the secure sign-in process.
    \item \textbf{Asynchronous AI Execution Pipeline:} Figure~\ref{fig:ai_architecture} in Chapter~\ref{chap:ai-design} (Page~\pageref{fig:ai_architecture}) details prompt formatting and structured JSON validations.
    \item \textbf{Zero-Egress Direct Storage Flow:} Figure~\ref{fig:storage_flow_diagram} in Chapter~\ref{chap:low-level-design} (Page~\pageref{fig:storage_flow_diagram}) shows the pre-signed URL upload lifecycle.
\end{itemize}

\section{Key Architecture Statistics}
Table~\ref{tab:architecture_statistics} provides the structural dimensions of the Nexus PM platform as of the current production release.

\begin{table}[h!]
\centering
\small
\renewcommand{\arraystretch}{1.4}
\begin{tabularx}{\textwidth}{>{\bfseries\color{primary}}l l X}
    \toprule[1.5pt]
    \rowcolor{primary!5}
    \textbf{\color{primary}Architectural Domain} & \textbf{\color{primary}Metric Value} & \textbf{\color{primary}Technical Details} \\
    \midrule
    Database Schemas & 11 Tables & Mapped via SQLAlchemy ORM models, tracked via Alembic versions. \\
    API Routes & 29 Endpoints & Partitioned across 9 routers (Auth, Projects, Members, Tasks, etc.). \\
    Security Protocols & 2 Identity Providers & Native credential signup/login and Google OAuth 2.0 federation. \\
    Storage Containers & 2 Locations & Supabase PostgreSQL (Structured) and Cloudflare R2 (Unstructured). \\
    Caching Providers & 1 Serverless Node & Upstash Redis instance handling sessions and rate-limits. \\
    Testing Level & 5 Build Quality Gates & Formatting, Linting, Dependency scan, Migration run, Unit tests. \\
    \bottomrule[1.5pt]
\end{tabularx}
\caption{Nexus PM Platform Architectural Dimensions}
\label{tab:architecture_statistics}
\end{table}

\section{Chapter Summary}
This chapter presented a high-level summary of the Nexus PM architecture, establishing the technologies, principles, deployment structures, and key statistics of the system. The next chapter will provide the background context of the platform, detailing its business drivers, target personas, system scope, and constraints in the Project Overview.
